\chapter{Iteración 4}

\begin{quotation}[Empresario]{Steve Jobs (1955--2011)}
El diseño no es solo lo que ves, sino cómo funciona.
\end{quotation}

\begin{abstract}
En este capítulo reflejaremos todos y cada uno de los costes que se han de reflejar a la hora de desarrollar el proyecto así como la relación de funcionalidades que se desarrollan con estos costes.
\end{abstract}

\section{Características a desarrollar}

\begin{enumerate}
\item Funcionalidad 1: Obtención de los datos de usuario. Ver Tabla \ref{tab:valor1AportadoIteracion4}.
\item Funcionalidad 2: Visualización de equipos. Ver Tabla \ref{tab:valor2AportadoIteracion4}.
\item Funcionalidad 3: Perfil del usuario. Ver Tabla \ref{tab:valor3AportadoIteracion4}.
\item Funcionalidad 4: Creación de invitaciones a equipos. Ver Tabla \ref{tab:valor4AportadoIteracion4}.
\item Funcionalidad 5: Aceptación y Rechazo de invitaciones. Ver Tabla \ref{tab:valor5AportadoIteracion4}.
\item Funcionalidad 6: Adición de persistencia. Ver Tabla \ref{tab:valor6AportadoIteracion4}.
\end{enumerate}

%----------------------------------------------------------------------------------------------------------

\begin{table*}[htb]
	\centering
	\begin{coolTable}{p{4cm}p{\textwidth-4.5cm}}{2}
{Análisis de valor aportado 0001}
\textbf{Propuesta}&Se pretende desarrollar la funcionalidad de obtener los datos de los usuarios registrados en la aplicación\\
\midrule
\textbf{Valor}&El valor aportado es tener un registro de los datos de los usuarios registrados.\\
\textbf{Coste}&Se requiere la participación de ambos miembros del equipo desarrollador para llevar a cabo la implementación.\\
\textbf{Opciones}&Esta funcionalidad tiene un valor alto en la aplicación.\\
\textbf{Riesgos}&La no correcta obtención de todos los datos del usuario.\\
\textbf{Deuda técnica}&Se ha elegido la solución más factible para el desarrollo de la misma.\\
	\end{coolTable}
	\caption{Análisis de valor aportado 0001\label{tab:valor1AportadoIteracion4}}
\end{table*}

%----------------------------------------------------------------------------------------------------------

\begin{table*}[htb]
	\centering
	\begin{coolTable}{p{4cm}p{\textwidth-4.5cm}}{2}
{Análisis de valor aportado 0002}
\textbf{Propuesta}&Se pretende desarrollar la funcionalidad de visualizar los equipos creados por los usuarios.\\
\midrule
\textbf{Valor}&El valor aportado es poder visualizar los equipos creados por los usuarios.\\
\textbf{Coste}&Se requiere la participación de ambos miembros del equipo desarrollador para llevar a cabo la implementación.\\
\textbf{Opciones}&Esta funcionalidad tiene un valor prioritario en la aplicación.\\
\textbf{Riesgos}&La posibilidad de no poder visualizar los equipos de forma satisfactoria.\\
\textbf{Deuda técnica}&Se ha elegido la solución más factible para el desarrollo de la misma.\\
	\end{coolTable}
	\caption{Análisis de valor aportado 0002\label{tab:valor2AportadoIteracion4}}
\end{table*}

%----------------------------------------------------------------------------------------------------------

\begin{table*}[htb]
	\centering
	\begin{coolTable}{p{4cm}p{\textwidth-4.5cm}}{2}
{Análisis de valor aportado 0003}
\textbf{Propuesta}&Se pretende desarrollar la funcionalidad del perfil de usuario.\\
\midrule
\textbf{Valor}&El valor aportado es poder visualizar los datos del usuario registrado.\\
\textbf{Coste}&Se requiere la participación de ambos miembros del equipo desarrollador para llevar a cabo la implementación.\\
\textbf{Opciones}&Esta funcionalidad tiene un valor alto en la aplicación.\\
\textbf{Riesgos}&La no correcta visualización de los diferentes datos que tiene el usuario.\\
\textbf{Deuda técnica}&Se ha elegido la solución más factible para el desarrollo de la misma.\\
	\end{coolTable}
	\caption{Análisis de valor aportado 0003\label{tab:valor3AportadoIteracion4}}
\end{table*}

%----------------------------------------------------------------------------------------------------------

\begin{table*}[htb]
	\centering
	\begin{coolTable}{p{4cm}p{\textwidth-4.5cm}}{2}
{Análisis de valor aportado 0004}
\textbf{Propuesta}&Se pretende desarrollar la funcionalidad de crear invitaciones para invitar a los usuarios al equipo.\\
\midrule
\textbf{Valor}&El valor aportado es poder invitar a los usuarios al equipo.\\
\textbf{Coste}&Se requiere la participación de ambos miembros del equipo desarrollador para llevar a cabo la implementación.\\
\textbf{Opciones}&Esta funcionalidad tiene un valor alto en la aplicación.\\
\textbf{Riesgos}&La posibilidad de no poder invitar a los usuarios al equipo.\\
\textbf{Deuda técnica}&Se ha elegido la solución más factible para el desarrollo de la misma.\\
	\end{coolTable}
	\caption{Análisis de valor aportado 0004\label{tab:valor4AportadoIteracion4}}
\end{table*}

%----------------------------------------------------------------------------------------------------------

\begin{table*}[htb]
	\centering
	\begin{coolTable}{p{4cm}p{\textwidth-4.5cm}}{2}
{Análisis de valor aportado 0005}
\textbf{Propuesta}&Se pretende desarrollar la funcionalidad de aceptar y rechazar las invitaciones.\\
\midrule
\textbf{Valor}&El valor aportado es decidir si aceptar una invitación o rechazarla si no es del equipo correcto.\\
\textbf{Coste}&Se requiere la participación de ambos miembros del equipo desarrollador para llevar a cabo la implementación.\\
\textbf{Opciones}&Esta funcionalidad tiene un valor alto en la aplicación.\\
\textbf{Riesgos}&La no correcta aceptación o rechazo de las invitaciones.\\
\textbf{Deuda técnica}&Se ha elegido la solución más factible para el desarrollo de la misma.\\
	\end{coolTable}
	\caption{Análisis de valor aportado 0005\label{tab:valor5AportadoIteracion4}}
\end{table*}

%----------------------------------------------------------------------------------------------------------------------------------------------------

\begin{table*}[htb]
	\centering
	\begin{coolTable}{p{4cm}p{\textwidth-4.5cm}}{2}
{Análisis de valor aportado 0006}
\textbf{Propuesta}&Se pretende desarrollar la funcionalidad de que los datos del usuario sean persistentes en la aplicación.\\
\midrule
\textbf{Valor}&El valor aportado es mantener la sesión del usuario que esta conectado en la aplicación.\\
\textbf{Coste}&Se requiere la participación de ambos miembros del equipo desarrollador para llevar a cabo la implementación.\\
\textbf{Opciones}&Esta funcionalidad tiene un valor prioritario en la aplicación.\\
\textbf{Riesgos}&La falta de persistencia del usuario en la aplicación.\\
\textbf{Deuda técnica}&Se ha elegido la solución más factible para el desarrollo de la misma.\\
	\end{coolTable}
	\caption{Análisis de valor aportado 0006\label{tab:valor6AportadoIteracion4}}
\end{table*}

%----------------------------------------------------------------------------------------------------------------------------------------------------

\clearpage

\section{Diseño}


Esto es comparable en la siguiente figura. \ref{fig:diseno04}.

\begin{figure}[htbp]
\begin{center}

\includegraphics[width=\textwidth]{figures/UML-iteracion-4.jpg}
\caption{Diagrama UML de diseño para la iteración 4}
\label{fig:diseno04}
\end{center}
\end{figure}

\clearpage
Un memorando técnico por cada decisión de diseño.

\begin{table*}[htb]
	\centering
	\begin{coolTable}{p{4cm}p{\textwidth-4.5cm}}{2}
{Memorando técnico 0001}
\textbf{Asunto}&Diseño del modelo Tarea, añadiendo fecha estimada de finalización y fecha real de finalización.\\
\textbf{Resumen}&Se ha diseñado un modelo de Tarea que contenga los elementos necesarios para nuestra definición de esta; título, descripción, fechas de inicio y fin, duración, priorización si está terminada o no, además contaremos con el creador de la tarea, el usuario asignado a la tarea y el equipo de la tarea si esta tuviera. Durante esta iteración se ha añadido la estimación de la finalización de la tarea.\\
\midrule
\textbf{Factores causantes}&La relación entre modelos, ya que la clase Tarea cuenta con una asociación tanto con el modelo Equipo como con el modelo Usuario, así como los datos que requerimos que estos tengan asignados para el proyecto. \\
\textbf{Solución}&El modelo resultante asigna un tipo de dato para cada uno de los campos así como las relaciones que no pueden ser nulas en ningún momento por la concepción del modelo de Tarea como pueden ser el nombre o el creador.\\
\textbf{Motivación}&Un modelo que contenga los datos necesarios para un correcto procesamiento de las tareas.\\
\textbf{Cuestiones abiertas}& La correcta relación entre tablas y completitud de los atributos.\\
\textbf{Alternativas}&Se valoraron otras alternativas de relaciones, pero la dispuesta es la que se consideró más óptima tanto a niveles de implementación como de seguridad además de ser óptimas en el tipo de cada variable del modelo.\\
	\end{coolTable}
	\caption{Memorando técnico 0001}
\end{table*}

%-----------------------------------------------------------------------------------------------------------
\begin{table*}[htb]
	\centering
	\begin{coolTable}{p{4cm}p{\textwidth-4.5cm}}{2}
{Memorando técnico 0002}
\textbf{Asunto}&Diseño del modelo Equipo.\\
\textbf{Resumen}&Se ha diseñado un modelo de Equipo que contenga los elementos necesarios para nuestra definición de esta; nombre, además contaremos con el creador del equipo, lista de asociados, así como de la lista de managers.\\
\midrule
\textbf{Factores causantes}&La relación entre modelos, ya que la clase Equipo cuenta con una asociación con el modelo Usuario, así como los datos que requerimos que estos tengan asignados para el proyecto.\\
\textbf{Solución}&El modelo resultante asigna un tipo de dato para cada uno de los campos.\\
\textbf{Motivación}&Un modelo que contenga los datos necesarios para un correcto procesamiento de los equipos.\\
\textbf{Cuestiones abiertas}& La correcta relación entre tablas y completitud de los atributos.\\
\textbf{Alternativas}&Dado a la tecnología utilizada la manera de llevar a cabo esta relación entre tablas es la más óptima.\\
	\end{coolTable}
	\caption{Memorando técnico 0002}
\end{table*}

%-------------------------------------------------------------------------------------------------------

\begin{table*}[htb]
	\centering
	\begin{coolTable}{p{4cm}p{\textwidth-4.5cm}}{2}
{Memorando técnico 0003}
\textbf{Asunto}&Diseño del modelo Usuario.\\
\textbf{Resumen}&Se ha diseñado un modelo de Usuario que contenga los elementos necesarios para nuestra definición de este; nombre, apellidos, correo electrónico, la hora de inicio de su jornada laboral, hora de fin de la jornada laboral, nombre de usuario, contraseña, una lista de equipos y un atributo boolean que denota si este ha iniciado sesión.\\
\midrule
\textbf{Factores causantes}&La relación entre modelos, ya que la clase Usuario cuenta con dos asociaciones con el modelo Tarea y con el modelo Equipo, así como los datos que requerimos que estos tengan asignados para el proyecto.\\
\textbf{Solución}&El modelo resultante asigna un tipo de dato para cada uno de los campos.\\
\textbf{Motivación}&Un modelo que contenga los datos necesarios para un correcto procesamiento de las Usuario.\\
\textbf{Cuestiones abiertas}&La correcta relación entre tablas y completitud de los atributos.\\
\textbf{Alternativas}&Se valoraron otras alternativas de relaciones, pero la dispuesta es la que se consideró más óptima tanto a niveles de implementación como de seguridad.\\
	\end{coolTable}
	\caption{Memorando técnico 0003}
\end{table*}

%----------------------------------------------------------------------------------------------------------
\begin{table*}[htb]
	\centering
	\begin{coolTable}{p{4cm}p{\textwidth-4.5cm}}{2}
{Memorando técnico 0004}
\textbf{Asunto}&Diseño del modelo Invitación.\\
\textbf{Resumen}&Se ha diseñado un modelo de Invitación que contenga los elementos necesarios para nuestra definición de esta; mensaje, estado, correo electrónico, además contaremos con el usuario invitado, así como el equipo al que pertenece la invitación.\\
\midrule
\textbf{Factores causantes}&La relación entre modelos, ya que la clase Invitación cuenta con una asociación con el modelo Usuario y con una asociación con el modelo Equipo, así como los datos que requerimos que estos tengan asignados para el proyecto.\\
\textbf{Solución}&El modelo resultante asigna un tipo de dato para cada uno de los campos.\\
\textbf{Motivación}&Un modelo que contenga los datos necesarios para un correcto procesamiento de las invitaciones.\\
\textbf{Cuestiones abiertas}& La correcta relación entre tablas y completitud de los atributos.\\
\textbf{Alternativas}&Dado a la tecnología utilizada la manera de llevar a cabo esta relación entre tablas es la más óptima.\\
	\end{coolTable}
	\caption{Memorando técnico 0004}
\end{table*}

%------------------------------------------------------------------------------------------------------

\clearpage

\section{Implementación}

En esta sección se va a realizar un memorando técnico por cada decisión de implementación y refactorización que afecte al diseño.

\begin{table*}[htb]
	\centering
	\begin{coolTable}{p{4cm}p{\textwidth-4.5cm}}{2}
{Memorando técnico 0001}
\textbf{Asunto}&La creación de la clase Tarea.\\
\textbf{Resumen}&Se han definido los siguientes atributos para la clase Tarea que son los necesarios para mostrar los datos que queremos.\\
\midrule
\textbf{Factores causantes}&La relación entre tablas, ya que la clase Tarea cuenta con una asociación tanto con la clase Equipo como con la clase Usuario. \\
\textbf{Solución}&La solución llevada a cabo tiene la utilización de las anotaciones que proporciona JPA, para gestionar las relaciones entre dos tablas. En este caso para la relación entre Tarea y Equipo será \textbf{@ManyToOne} y para la relación entre Tarea y Usuario será \textbf{@OneToOne}.\\
\textbf{Motivación}&Ya conocíamos esta forma de relacionar las tablas.\\
\textbf{Cuestiones abiertas}& La correcta relación entre tablas.\\
\textbf{Alternativas}&Dado a la tecnología utilizada la manera de llevar a cabo esta relación entre tablas es la más óptima.\\
	\end{coolTable}
	\caption{Memorando técnico 0001}
\end{table*}

%-----------------------------------------------------------------------------------------------------------
\begin{table*}[htb]
	\centering
	\begin{coolTable}{p{4cm}p{\textwidth-4.5cm}}{2}
{Memorando técnico 0002}
\textbf{Asunto}&La creación de la clase Equipo.\\
\textbf{Resumen}&Se han definido los siguientes atributos para la clase Equipo que son los necesarios para mostrar los datos que queremos.\\
\midrule
\textbf{Factores causantes}&La relación entre tablas, ya que la clase Equipo cuenta con una asociación con la clase Usuario. \\
\textbf{Solución}&La solución llevada a cabo tiene la utilización de las anotaciones que proporciona JPA, para gestionar las relaciones entre dos tablas. En este caso la relación entre Equipo y Usuario serán \textbf{@ManyToMany}.\\
\textbf{Motivación}&Ya conocíamos esta forma de relacionar las tablas.\\
\textbf{Cuestiones abiertas}& La correcta relación entre tablas.\\
\textbf{Alternativas}&Dado a la tecnología utilizada la manera de llevar a cabo esta relación entre tablas es la más óptima.\\
	\end{coolTable}
	\caption{Memorando técnico 0002}
\end{table*}

%-------------------------------------------------------------------------------------------------------

\begin{table*}[htb]
	\centering
	\begin{coolTable}{p{4cm}p{\textwidth-4.5cm}}{2}
{Memorando técnico 0003}
\textbf{Asunto}&La creación de la clase Usuario.\\
\textbf{Resumen}&Se han definido los siguientes atributos para la clase Usuario que son los necesarios para mostrar los datos que queremos.\\
\midrule
\textbf{Factores causantes}&La relación entre tablas, ya que la clase Usuario cuenta con una asociación con la clase Equipo. \\
\textbf{Solución}&La solución llevada a cabo tiene la utilización de las anotaciones que proporciona JPA, para gestionar las relaciones entre dos tablas. En este caso para la relación entre Usuario y Equipo será \textbf{@OneToMany}.\\
\textbf{Motivación}&Ya conocíamos esta forma de relacionar las tablas.\\
\textbf{Cuestiones abiertas}& La correcta relación entre tablas.\\
\textbf{Alternativas}&Dado a la tecnología utilizada la manera de llevar a cabo esta relación entre tablas es la más óptima.\\
	\end{coolTable}
	\caption{Memorando técnico 0003}
\end{table*}

%----------------------------------------------------------------------------------------------------------
\begin{table*}[htb]
	\centering
	\begin{coolTable}{p{4cm}p{\textwidth-4.5cm}}{2}
{Memorando técnico 0004}
\textbf{Asunto}&La creación de la clase Invitación.\\
\textbf{Resumen}&Se han definido los siguientes atributos para la clase Invitación que son los necesarios para mostrar los datos que queremos.\\
\midrule
\textbf{Factores causantes}&La relación entre tablas, ya que la clase Invitación cuenta con una asociación tanto con la clase Equipo como con la clase Usuario. \\
\textbf{Solución}&La solución llevada a cabo tiene la utilización de las anotaciones que proporciona JPA, para gestionar las relaciones entre dos tablas. En este caso para la relación entre Invitación y Equipo será \textbf{@OneToOne} y para la relación entre Invitación y Usuario será \textbf{@OneToOne}.\\
\textbf{Motivación}&Ya conocíamos esta forma de relacionar las tablas.\\
\textbf{Cuestiones abiertas}& La correcta relación entre tablas.\\
\textbf{Alternativas}&Dado a la tecnología utilizada la manera de llevar a cabo esta relación entre tablas es la más óptima.\\
	\end{coolTable}
	\caption{Memorando técnico 0004}
\end{table*}

%-----------------------------------------------------------------------------------------------------

\clearpage

\section{Pruebas}

De la misma manera que con la iteración anterior no se han integrado ninguna prueba automatizada sobre el código, estas se desarrollarán en su totalidad en la sexta iteración. Por lo que por ahora aunque no se han desarrollado pruebas automatizadas, el software producido ha sido testeado por los desarrolladores de manera manual.

\section{Despliegue}

Durante esta iteración se han realizado el mismo número de despliegues que en la iteración anterior es decir se han 5 despliegues semanales, los lunes, miércoles, viernes, sábados y domingos, sobre la plataforma de Heroku.